% ===== PRÉAMBULE CANONIQUE — à coller VERBATIM en tête de CHAQUE .tex =====
\documentclass[11pt,a4paper]{article}
\usepackage[utf8]{inputenc}\usepackage[T1]{fontenc}\usepackage{lmodern}
\usepackage[french]{babel}
\usepackage{amsmath,amssymb,mathtools}\usepackage{icomma}
\usepackage[version=4]{mhchem}          % \ce{...} : équations & formules
\usepackage{siunitx}\sisetup{locale=FR} % \qty{}{}, \si{}, \ang{}
\usepackage{chemfig}                    % \chemfig{...} : structures développées
\usepackage{xcolor}\usepackage{colortbl}
\usepackage{tikz}\usepackage{pgfplots}\pgfplotsset{compat=1.18}
\usepackage[most]{tcolorbox}\usepackage{enumitem}\usepackage{array,booktabs}
\usepackage[a4paper,margin=2.2cm]{geometry}
\usetikzlibrary{arrows.meta,calc,angles,quotes,positioning,patterns,backgrounds,shapes.geometric}
\definecolor{primary}{RGB}{222,49,99}\definecolor{primarydark}{RGB}{150,22,55}
\definecolor{defcol}{RGB}{0,114,178}\definecolor{methcol}{RGB}{52,92,148}
\definecolor{excol}{RGB}{0,135,90}\definecolor{attncol}{RGB}{200,40,55}
\tcbset{boxbase/.style={enhanced,breakable,boxrule=0.9pt,arc=2.5pt,
  left=3mm,right=3mm,top=2.3mm,bottom=2.3mm,fonttitle=\bfseries,coltitle=white}}
% --- boîtes TUTORIEL (noms EXACTS, ne pas renommer) ---
\newtcolorbox{cle}[1][]{boxbase,colback=defcol!6,colframe=defcol,
  title={\textsc{À retenir}\ifx\relax#1\relax\relax\else\textmd{ : \itshape #1}\fi}}
\newtcolorbox{methode}[1][]{boxbase,colback=methcol!7,colframe=methcol,sharp corners,
  title={\textsc{Méthode}\ifx\relax#1\relax\relax\else\textmd{ --- \itshape #1}\fi}}
\newtcolorbox{exemplebox}[1][]{boxbase,colback=excol!5,colframe=excol,
  title={\textsc{Exemple}\ifx\relax#1\relax\relax\else\textmd{ : \itshape #1}\fi}}
\newtcolorbox{piege}[1][]{boxbase,colback=attncol!6,colframe=attncol,
  title={\textsc{Piège}\ifx\relax#1\relax\relax\else\textmd{ : \itshape #1}\fi}}
\newtcolorbox{remarque}[1][]{boxbase,colback=black!4,colframe=black!45,
  title={\textsc{Remarque}\ifx\relax#1\relax\relax\else\textmd{ : \itshape #1}\fi}}
% --- boîtes EXERCICE / CORRIGÉ (noms EXACTS, auto-comptées) ---
\newtcolorbox[auto counter]{exo}[1][]{boxbase,colback=primary!4,colframe=primary,
  title={\textsc{Exercice}~\thetcbcounter\ifx\relax#1\relax\relax\else\textmd{ --- \itshape #1}\fi}}
\newtcolorbox[auto counter]{corrige}[1][]{boxbase,colback=excol!4,colframe=excol,
  title={\textsc{Corrigé}~\thetcbcounter\ifx\relax#1\relax\relax\else\textmd{ --- \itshape #1}\fi}}
\newcommand{\R}{\mathbb{R}}\newcommand{\N}{\mathbb{N}}\newcommand{\Z}{\mathbb{Z}}\newcommand{\ds}{\displaystyle}
% (un \usepackage{fancyhdr}+en-tête est possible ; il est ignoré côté web)
% =========================================================================
\begin{document}
\begin{corrige}[covalente ou ionique ?]
Deux atomes identiques ou deux non-métaux $\Rightarrow$ covalente ; un métal $+$ un non-métal
$\Rightarrow$ ionique.
\begin{enumerate}[label=\alph*)]
  \item \ce{Cl2} : \textbf{covalente} (deux atomes identiques, aucun n'est plus avide d'électrons).
  \item \ce{KCl} : \textbf{ionique} (\ce{K} métal $+$ \ce{Cl} non-métal, fort écart).
  \item \ce{O2} : \textbf{covalente} (deux atomes identiques).
  \item \ce{MgO} : \textbf{ionique} (\ce{Mg} métal $+$ \ce{O} non-métal).
  \item \ce{HCl} : \textbf{covalente} (deux non-métaux ; liaison polarisée mais non ionique).
\end{enumerate}
\end{corrige}

\begin{corrige}[électrons de valence d'un atome]
Le nombre d'électrons de valence est le numéro de colonne de l'atome.
\begin{enumerate}[label=\alph*)]
  \item \ce{N} : $\mathbf{5}$.
  \item \ce{O} : $\mathbf{6}$.
  \item \ce{C} : $\mathbf{4}$.
  \item \ce{Cl} : $\mathbf{7}$.
  \item \ce{H} : $\mathbf{1}$.
\end{enumerate}
\end{corrige}

\begin{corrige}[octet ou duet ?]
Seuls \ce{H} et \ce{He} visent le duet ($2$~e$^-$) ; tous les autres visent l'octet ($8$~e$^-$).
\begin{enumerate}[label=\alph*)]
  \item \ce{H} : \textbf{duet}.
  \item \ce{C} : \textbf{octet}.
  \item \ce{O} : \textbf{octet}.
  \item \ce{He} : \textbf{duet}.
  \item \ce{F} : \textbf{octet}.
\end{enumerate}
\end{corrige}

\begin{corrige}[l'octet est-il respecté ?]
On compte $2$~e$^-$ par liaison et $2$~e$^-$ par doublet non liant.
\begin{enumerate}[label=\alph*)]
  \item \ce{N} : $3\times 2 + 2 = 8$~e$^-$ $\Rightarrow$ \textbf{octet respecté}.
  \item \ce{C} : $4\times 2 = 8$~e$^-$ $\Rightarrow$ \textbf{octet respecté}.
  \item \ce{O} : $2\times 2 + 2\times 2 = 8$~e$^-$ $\Rightarrow$ \textbf{octet respecté}.
  \item \ce{H} : $1\times 2 = 2$~e$^-$ $\Rightarrow$ \textbf{duet respecté}.
\end{enumerate}
\end{corrige}

\begin{corrige}[doublets liants et non liants]
On lit directement la structure de Lewis.
\begin{enumerate}[label=\alph*)]
  \item $\mathbf{3}$ doublets liants (les trois liaisons \ce{N-H}).
  \item $\mathbf{1}$ doublet non liant (le lobe violet sur l'azote).
  \item Oui : $3\times 2 + 2 = 8$~e$^-$ autour de \ce{N}, l'\textbf{octet est respecté}.
\end{enumerate}
\end{corrige}

\end{document}
